\chapter{Discussion}
\label{chap:discussion}

\section{When Simple Schedules Are Near-Optimal}
\label{sec:disc:nearoptimal}

\TODO{Characterise the regime (symmetric noise, uniform hops, moderate
budget) in which the paper's hand-chosen YY-first schedule is within
a small fraction of the DP optimum. Provide an intuitive explanation
based on the error-vector bias toward ZI/IZ Pauli errors.}

\section{When the Scheduling Choice Matters}
\label{sec:disc:matters}

\TODO{Identify the regime (asymmetric per-hop noise, tight fidelity
constraint, large $N$) where the optimised schedule achieves a
meaningful improvement. Discuss the excluded move (end-to-end
purification of partially-purified multi-hop segments) and its impact.}

\section{Optimistic versus Heralded Purification}
\label{sec:disc:optimistic}

\TODO{Analyse the trade-off: optimistic purification saves one
classical round-trip ($L/c$ vs.\ $2L/c$ memory time) at the cost of
higher discard probability. Show which regime favours each mode.}

\section{Limitations and Threats to Validity}
\label{sec:disc:limitations}

\TODO{State the main modelling assumptions (Pauli channel, i.i.d.\ noise,
non-adaptive schedules, single-path routing) and discuss what could
change if they were relaxed.}